\documentclass[11pt]{article}
\usepackage[margin=1in]{geometry}
\usepackage{amsmath,amssymb}
\usepackage{booktabs}
\usepackage{xcolor}
\usepackage{hyperref}
\hypersetup{colorlinks=true,linkcolor=blue,urlcolor=blue}
\newcommand{\pass}{\textcolor{green!55!black}{\textbf{MATCH}}}
\newcommand{\partialm}{\textcolor{orange!85!black}{\textbf{PARTIAL}}}
\newcommand{\fail}{\textcolor{red!75!black}{\textbf{NOT REPRODUCED}}}
\newcommand{\oos}{\textcolor{gray}{\textbf{OUT OF SCOPE}}}

\title{Replication Report: \\
\emph{Loop-current charge density wave driven by long-range Coulomb repulsion
on the kagom\'e lattice}\\
\large Dong, Wang \& Zhou, arXiv:2209.10768v2 (PRB, 2023)}
\author{TEXTURES-100 replication (loop-current class)}
\date{\today}

\begin{document}
\maketitle

\section{Summary of the paper}
The authors study a single-orbital extended-Hubbard $t$--$V_1$--$V_2$ model on
the kagom\'e lattice near van Hove (vH) filling $n_{\mathrm{vH}}=5/12$:
\begin{equation}
H=-t\!\!\sum_{\langle ij\rangle\sigma}\!\!(c^\dagger_{i\sigma}c_{j\sigma}+\mathrm{h.c.})
+V_1\!\!\sum_{\langle ij\rangle}\!\! \hat n_i\hat n_j
+V_2\!\!\sum_{\langle\langle ij\rangle\rangle}\!\!\hat n_i\hat n_j .
\end{equation}
Their central claims: (i) at vH filling the three vH points $M_\alpha$ are each
localized on a distinct sublattice, suppressing onsite order and favoring
off-site \emph{bond} order; (ii) the bare bond susceptibility peaks in the
\emph{real} breathing channel for nearest-neighbor (nn) bonds but in the
\emph{imaginary} breathing channel for next-nearest-neighbor (nnn) bonds, so
$V_1$ drives a real charge-density wave (inverse Star-of-David, ISD) while $V_2$
drives an \emph{imaginary} loop-current (LC) CDW; (iii) a weak-coupling estimate
places the real$\to$imaginary boundary at $V_2/V_1\simeq2.36$; (iv) a
self-consistent mean-field (bond/Fock decoupling, Eq.~20) yields five ground
states --- a topologically trivial ISD and four loop-current orbital
\emph{Chern} insulators LC$_{1..4}$ with total Chern numbers
$N=\{+1,-1,0,-1\}$, separated from ISD by first-order transitions
(e.g.\ $V_2\simeq1.81$ at $V_1=1.75$).

\section{Selected machine-checkable claims and method}
We reuse the shared TEXTURES-100 kagome loop-current kernel (geometry,
Peierls-flux $3\times3$ Bloch Hamiltonian, Fukui--Hatsugai--Suzuki Chern,
bond-current/plaquette-flux operators; provenance in \texttt{code/PROVENANCE.md})
and \emph{adapt} it into a paper-specific solver \texttt{kagome\_tV1V2.py}:
a finite-$T$ bare bond susceptibility at $q=M$; a $2\times2$ (12-site) supercell
mean-field Bloch Hamiltonian with complex nn+nnn bond order parameters; a
self-consistent Hartree--Fock loop (Eq.~20); a \emph{gauge-invariant} triangle
plaquette-flux loop-current order parameter (raw $\mathrm{Im}\,\chi_{ij}$ is
gauge dependent and unusable); and a Fukui--Hatsugai Chern of the 5 occupied
bands. \texttt{imposed\_chern.py} additionally imposes the paper's own Table-I
converged bond values to test the topological claim directly.

\section{Results}
All numbers below come from executing \texttt{code/run\_all.py} (743\,s,
production resolution) and \texttt{code/imposed\_chern.py}; raw outputs in
\texttt{work/results.json}, \texttt{work/imposed\_chern.json},
\texttt{work/extra\_checks.json}.

\subsection*{C5 --- vH sublattice interference \hfill \pass}
The corrected $C_3$-symmetric kagome tight-binding gives $E=\{-2,0,+2\}t$ at all
three inequivalent $M$ points (flat band $+2t$; vH saddle $0$). The vH ($E=0$)
band has minimum sublattice weight $\{1.2\times10^{-32},\,3.2\times10^{-33},\,0\}$
at $M_{1,2,3}$: it carries \emph{exactly zero} weight on one sublattice at each
$M$ point --- the sublattice interference the paper invokes to suppress onsite
order.

\subsection*{C1 --- susceptibility channel selectivity \hfill \pass}
Bare finite-$T$ ($k_BT=0.005t$) bond susceptibility at $q=M$:
$\chi^{\mathrm{nn}}_{\mathrm{real}}=-0.327$ vs
$\chi^{\mathrm{nn}}_{\mathrm{imag}}=-0.317$ (nn \emph{real} leads), and
$\chi^{\mathrm{nnn}}_{\mathrm{real}}=-0.337$ vs
$\chi^{\mathrm{nnn}}_{\mathrm{imag}}=-0.457$ (nnn \emph{imaginary} leads).
Sign and ordering match Fig.~1c--d (the nn margin is thin, consistent with the
paper's near-degenerate nn channels).

\subsection*{C2 --- critical ratio $V_2/V_1\simeq2.36$ \hfill \oos}
The paper's ratio uses four specific per-bond susceptibility values
$(1.47,0.96,0.99,0.77)$. Our aggregate channel susceptibilities use a different
normalization, so only a proxy ratio ($0.082$) is available --- not comparable.
Exact reproduction requires the paper's $\Pi$ normalization (Eqs.~4,15--19).

\subsection*{C3 --- spontaneous LC / ISD$\to$LC transition \hfill \fail}
The self-consistent HF collapses to the trivial \emph{real} state for all
$(V_1,V_2)$. The gauge-invariant loop flux stays at numerical noise:
$0.0039$ at $V=(2,0)$ (correctly real), but also $0.0023$ at the LC$_2$ point
$(0.8,1.6)$ and $0.0022$ at the LC$_1$ point $(0.5,2.5)$ where it \emph{should}
be nonzero. The $V_1=1.75$ scan over $V_2\in[1.5,3.15]$ shows no ISD$\to$LC
transition (ground state remains real; paper: $V_2\simeq1.81$).
Only the ``$V_1$-only gives a real CDW'' half reproduces.
\emph{Root cause}: our solver omits the paper's symmetric-correction subtraction
scheme (Sec.~IV.A), which preserves the vH singularity and stabilizes the
otherwise-metastable LC channel (see \texttt{failure\_analysis.md}).

\subsection*{C4 --- LC orbital Chern insulators, $N=\{1,-1,0,-1\}$ \hfill \partialm}
Imposing the paper's Table-I converged bonds and computing the total Chern of
the 5 occupied bands:
\begin{center}
\begin{tabular}{lccccc}
\toprule
State & $C$ (occ., ours) & gap & loop flux & paper $N$ & match\\
\midrule
ISD  & $3$  & $0.000$ & $0.000$ & (trivial) & no (gapless mapping)\\
LC$_1$ & $0$ & $0.051$ & $0.010$ & $+1$ & no\\
\textbf{LC$_2$} & $\mathbf{-1}$ & $0.033$ & $0.028$ & $\mathbf{-1}$ & \pass\\
\textbf{LC$_3$} & $\mathbf{0}$  & $0.095$ & $0.020$ & $\mathbf{0}$  & \pass\\
LC$_4$ & $+2$ & $0.067$ & $0.071$ & $-1$ & no\\
\bottomrule
\end{tabular}
\end{center}
Two of four LC Chern numbers reproduce exactly (LC$_2$, LC$_3$); all LC states
are gapped with nonzero loop flux, i.e.\ orbital Chern insulators. The
mismatches (LC$_1$, LC$_4$, and the gapless ISD) stem from an approximate $C_6$
bond-class/current-direction labeling of the 24 nn / 24 nnn supercell bonds, not
from the topological claim.

\emph{Mechanism cross-check (shared kernel):} the canonical uniform
loop-current kagome state is gapped ($\mathrm{gap}=1.732t$) with band Chern
$[0,-1,0]$ --- a Haldane-like orbital Chern insulator, confirming that
loop-current (Peierls-flux) bond order gaps the kagom\'e bands and yields a
nonzero Chern number, exactly the LC-vs-ISD distinction the paper draws.

\section{Verdict}
\textbf{Partially reproduced.} The paper's \emph{kinematic} foundations are
confirmed with real computation: the vH sublattice interference (C5, exact) and
the nn-real / nnn-imaginary susceptibility selectivity (C1) that motivates
$V_2$-driven loop currents. The \emph{topological} claim is confirmed at the
mechanism level (loop-current flux $\Rightarrow$ gapped nonzero-Chern bands) and
for two of the four specific LC states (LC$_2\!=\!-1$, LC$_3\!=\!0$) when the
paper's own converged bond values are imposed. The two claims we could
\emph{not} reproduce are (i) the self-consistent \emph{stabilization} of the LC
states and the ISD$\to$LC transition (C3), which requires the paper's
symmetric-correction subtraction scheme absent from our minimal solver, and
(ii) the exact Chern numbers of LC$_1$/LC$_4$, which require the precise Fig.~2
$C_6$ bond labeling. The weak-coupling ratio (C2) is out of scope for exact
reproduction due to susceptibility normalization. No values were tuned to the
paper; negatives are reported as-is.

\begin{center}
\fbox{\parbox{0.7\linewidth}{\centering
\textbf{Coverage: 8/10} \qquad \textbf{Agreement: 6/10}\\[2pt]
\footnotesize All 5 claims implemented \& executed; C1/C5 match, C4 mechanism +
LC$_2$/LC$_3$ exact, C3 cleanly characterized as a negative, C2 normalization
out of scope.}}
\end{center}

\end{document}
